\documentclass[11pt]{article}
\usepackage{fontspec}
\setmainfont{Times New Roman}
\setsansfont{Arial}
\setmonofont{Consolas}
\usepackage{amsmath,amssymb}
\usepackage{geometry}
\usepackage{microtype}
\geometry{a4paper,margin=3cm}
\setlength{\parindent}{1.25em}
\setlength{\parskip}{0pt}
\emergencystretch=30em
\allowdisplaybreaks
\raggedbottom
\providecommand{\mM}{\mathfrak M}
\providecommand{\mN}{\mathfrak N}
\providecommand{\mA}{\mathfrak A}
\providecommand{\mB}{\mathfrak B}
\providecommand{\mP}{\mathfrak P}
\providecommand{\mD}{\mathfrak D}
\providecommand{\mS}{\mathfrak S}
\providecommand{\ma}{\mathfrak a}
\providecommand{\mb}{\mathfrak b}

\begin{document}

% Унит: Папер 17 Сецтион 09. Дирецт Интерславиц Латин транслатион фром те Герман финал-аудитед соурце слице.

\setcounter{footnote}{16}
\subsection*{\S~9. Свез меджу појмами ``изоморфны'' и ``једнакого вида''.}

В тутом параграфу показујемо, же појем ``једнакого вида'', знаны при диференциалных изразах једнеј пременној, ако јест смыслно обобчен, води к изоморфији належечих груп остатков.

\textbf{Дефиниција.}\footnote{За диференциалне изразы једнеј пременној сравни формалну дефиницију Блумберга, на означеном месту, с. 25.} \textit{Два модула \(\mM\) и \(\mN\) называјут се модулами једнакого вида, ако обставајут два полиномы \(P\) и \(Q\), при чем \(P\) јест взаимно просты с \(\mM\), а \(Q\) взаимно просты с \(\mN\), тако же за всако \(M\equiv0(\mM)\) и \(N\equiv0(\mN)\)}
\begin{align}
 MQ&\equiv0(\mN),\tag{25}\\
 NP&\equiv0(\mM)\tag{26}
\end{align}
\textit{важи. Взаимна простост тут значи обставанье двох полиномов \(X\) и \(Y\), тако же}
\begin{align}
 YQ&\equiv1(\mN),\tag{27}\\
 XP&\equiv1(\mM)\tag{28}
\end{align}
\textit{важи.}

Ныне важи следујучи

\textbf{Теорем VI.} \textit{Два модула сут тогда и само тогда једнакого вида, когда јих групы остатков сут изоморфне.}

1. Неха \(\mA\) буде група остатков модула \(\mM\), \(\mB\) група остатков модула \(\mN\), и неха \(\mA\) јест изоморфна к \(\mB\); \(\ma\) и \(\mb\) неха будут јединичне класы од \(\mA\) и \(\mB\). Далеј неха \(Q\mb\) буде клас в \(\mB\), кторы одповеда јединици \(\ma\) в \(\mA\), и подобно \(P\ma\) клас в \(\mA\), кторы одповеда јединици \(\mb\) од \(\mB\). То придруженье записујемо в форме
\begin{align}
 \ma&\sim Q\mb,\tag{29}\\
 P\ma&\sim\mb.\tag{30}
\end{align}
Из того следује
\[
 0=M\ma\sim MQ\mb,
 \qquad\text{то јест}\qquad MQ\equiv0(\mN),
\]
чим (25) јест доказано. Такоже важи
\[
 NP\ma\sim N\mb=0,
 \qquad\text{то јест}\qquad NP\equiv0(\mM),
\]
теды релација (26). Далеј по (29) јест \(P\ma\sim PQ\mb\), и понеже по (30) было предположено \(P\ma\sim\mb\), из једнозначности придруженьа следује \(PQ\mb=\mb\), теды \(PQ\equiv1(\mN)\) и одповедајуче \(QP\equiv1(\mM)\). Тым (27) и (28) сут доказане с \(X=Q\), \(Y=P\).

2. Неха, обратно, \(\mM\) и \(\mN\) сут једнакого вида, тако же релације (25)--(28) обставајут. Тогда либовольному класу \(F\ma\) из \(\mA\) придружујемо клас \(FQ\mb\) из \(\mB\). При том всакому класу в \(\mA\) одповеда само једин клас в \(\mB\); бо из \(F\ma=0\), теды \(F=M\), по (25) следује \(FQ\mb=0\). Далеј цела група \(\mB\) јест изчрпана тојим придруженьем; бо специалному класу \(Y\ma\) одповеда клас \(YQ\mb\), кторы по (27) јест равны \(\mb\). Теды имајемо једнозначно определьено пресликанье \(\Phi\) из \(\mA\) на \(\mB\), але јошче не јест показано, же оно јест обратно једнозначно. Такоже из једначб (26) и (28) следује једнозначно определьено пресликанье \(\Psi\) из \(\mB\) на \(\mA\), кторо так само јошче не јест доказано како обратно једнозначно.

Ако једно из тих двох пресликань јест обратно једнозначно, тогда изоморфија јест доказана, понеже потребованьа гледе сумы и продукта сут изполньене. Но ако бы оба пресликаньа \(\Phi\) и \(\Psi\) не была обратно једнозначне, тогда бы обставали, на примјер в \(\mA\), класы \(G\ma\), кторым в \(\mB\) одповеда клас \(GQ\mb=0\). Класы \(FQ\mb\), кторе одповедајут осталим класам остатков \(F\ma\), тогда бы уже изчрпавали групу \(\mB\), и последовательно бы се через пресликанье \(\Psi\) в класах \(FQP\ma\) возвратила цела група \(\mA\). Како знајемо, обставајут такоже в \(\mB\) од нуле различне класы, кторым по \(\Psi\) одповеда нулевы клас в \(\mA\); означајемо с \(G_1\ma\) все те класы, за кторе \(G_1\ma\ne0\), але \(G_1QP\ma=0\). Вси остали полиномы \(F\) имајут својство, же \(FQP\ma\) изчрпава групу \(\mA\). Те полиномы \(F\), за кторе \(FQP\ma=G_1\ma\), означајемо с \(G_2\); таке полиномы обставајут всакы раз, когда обставајут полиномы \(G_1\). За всако \(G_2\) далеј важи \(G_2QP\ma\ne0\), зато \(G_2(QP)^2\ma=0\). Одповедајуче за всако цело число \(\nu\) обставајут класы \(G_\nu\ma\), тако же \(G_\nu(QP)^{\nu-1}\ma\ne0\), але \(G_\nu(QP)^\nu\ma=0\).

Ныне разматрјамо систему \(\mS\), составјену из всех полиномов \(G_1,G_2,\ldots\). Она има конечны модулны базис \(G_{\lambda_1},\ldots,G_{\lambda_e}\). Ныне выбира се \(\nu\) вече од највечего из индексов \(\lambda_1,\ldots,\lambda_e\), кторы означајемо с \(\lambda\). Тогда
\[
 G_\nu=A_1G_{\lambda_1}+\cdots+A_eG_{\lambda_e}.
\]
Ако то помножимо с \((QP)^\lambda\), доставајемо
\[
 G_\nu(QP)^\lambda\equiv0(\mM),
\]
что дава противреченье нашему предположеньу. Следује теды изоморфија, и зато по 1. важе острејше релације (27), (28) с \(X=Q\), \(Y=P\).

По Теорему VI можно Теорем V так само выразити в следујучеј форме:

\textbf{Теорем VII.}\footnote{Сравни Лоевы, с. 100--101.} \textit{Ако пополно редуцибилны модул \(\mM\) јест на два способы представимы како најмене обче многократно простых модулов \(\mP_1,\ldots,\mP_k\) односно \(\mD_1,\ldots,\mD_l\), кторе всакы раз творе тотално взаимно просту систему, тогда всакы \(\mP\) јест једнакого вида најмене с једним \(\mD\), и обратно. Ако всакы \(\mP\) јест једнакого вида точно с једним \(\mD\), тогда разклады сут идентичне. В противном случају најмене два из модулов \(\mP\) сут меджу собоју једнакого вида, и так само најмене два из модулов \(\mD\) сут меджу собоју једнакого вида.}

\end{document}

